%! Author = chtompki
%! Date = 1/14/26
% Example LaTeX document for GP111 - note % sign indicates a comment
\documentclass[12pt]{amsart}
% Default margins are too wide all the way around. I reset them here
\setlength{\hoffset}{-.75in}
\setlength{\textwidth}{6.5in}
\setlength{\voffset}{-.5in}
\setlength{\textheight}{9.0in}
\usepackage{amsmath}
\usepackage{amssymb}
\usepackage{amsthm}
\usepackage{pgfplots}
\usepackage{tikz}
\usetikzlibrary{arrows.meta} % Required for the Circle and Triangle arrow tips
\usepackage{setspace}
\title{Seminar Summary:\\January 27, 2026 - Rebecca Segal\\Numerical Solutions of Biologically Based Problems using Differential Equations}
\author{Rob Tompkins}
\date{\today}


\begin{document}
    \maketitle
    \date{\today}
    \begin{onehalfspace}
        In general across all the projects we explore below in this summary we have a general
        framework around how one of these projects is structured.
        Specifically the goal is to collaborate with biologist to understand research questions
        and the current state of the field.
        From this collaboration, our goal is to develop mathematical models that accommodate
        using numerical techniques for solving them.
        We use the biological data to paramaterize the model and/or run a sweep of the parameter
        space.
        Though it is worth noting that there is a natural trade-off between mathematical
        complexity of the model and the solvability of the model.

        The projects that we examined specifically were:
        \begin{itemize}
            \item Wound healing and inflamation
            \item Wound healing in coral
            \item Gut health and inflammation
            \item Nematocyst firing
            \item Phage therapy
            \item Sickle Cell Pain Episode Prediction
        \end{itemize}
        Our summaries on projects 5. and 6. were generally short because those both were described i
        n the last five-minutes of class.

        \section{Wound healing and inflammation.} For this project the original data came from
        an experiment that the original biologist did on himself (with medical assistance).
        Spefically he had a ``PTFE'' (or vortex) tube inserted under his skin for a period
        of time to create the wound in question as a mechanism for creating a standardized wound.
        Not surprisingly, there were few other participants in the clinical trial.
        That said, because it was biologist himself doing the experiment the richness of the data
        is interesting.
        For example, that there were even mounted slides from the experiment that showed the
        PTFE tube matrix along with fibroblasts laying down collagen.

        The model in question was a partial differential equation model based on the variables
        $m$ for the density of damaged the tissue, $n$ for the activated inflammatory cells, $f$ for the
        fibroblasts, and $c$ for new collagen deposition.
        We have four model equations, starting with the first for phagocytosis:
        \begin{displaymath}
            \frac{\partial m}{\partial t} = -\mu_{mn} mn
        \end{displaymath}
        The second equation represented diffusion and chemotaxis and deactivation, given by:
        \begin{displaymath}
            \frac{\partial n}{\partial t} = \nabla \cdot\left[D_n \nabla n - n \chi_n\nabla m\right] - \mu_{mn}mn
        \end{displaymath}
        The third equation represented diffustion/motility and chemotaxis and cell division and cell death
        \begin{displaymath}
            \frac{\partial f}{\partial t} = \nabla \cdot\left[D_f\nabla f - f\chi_f\nabla n\right] + s_f(n)f\left(2 - \frac{f}{F_0}\right) - \mu_f f
        \end{displaymath}
        And the fourth model equation represents deposition by $f$ and collateral damage
        \begin{displaymath}
            \frac{\partial c}{\partial t} = k_c f(C_0 - c) - \mu_{cn}cn.
        \end{displaymath}
        We were then presented some plots that represented collagen and fibroblast levels versus days
        with modeled data and real data ploted on the graphs.

        \section{Wound healing in coral.}
        The goal in this project is to consider coral in place of human tissue.
        Biologically we know that the damage processes (laceration, avulsion, fracture),
        inflammation processes, and physiological process for tissue regrowth behave in
        a similar fashion to humans.
        We noted that it was a convenient wound healing mechanism to study because the results
        could indeed be produced with coral in the lab as well as naturally occurring coral.
        We ask questions like whether the wound is regrowing or propagating.
        There were two specefied aims of the project, namely cellular level-wound healing
        versus tissue level-enegery focused

        In both of the aims of the project we began examining whether coral bleaching had any
        effect on the speed and mechanics of wound healing.
        Though it is know that glycerol production as a photosyntehtic byproduct transferred to
        mucus from the symbiotic algae yielded energy transfer to the calioblastic cells for
        skeletal growth.

        Potential future research projects out of this body of work could be examining
        multiple stressors, other coral species and linked cellular and tissue level models.

        \section{Gut health and inflammation.}
        So it's known that gut health has a considerable amount to do with body health from
        general inflammation, to mental health, to even levels of atherosclerosis.
        Now gut inflammation is regulated by an enzyme called alkaline phosphate (specifically
        intestinal alkaline phosphate or IAP).
        And, the gram-negative bacteria in the colon produce lipopolysaccharides (LPS) which are
        inflammatory agents.
        So the interplay between IAP and LPS regulates how the gut performs and even the structural
        integrity of wall of the bowel.

        If I understand the slides correctly our goal here is to write a mathematical model that
        effectively encompases the metabolic pathway for the breakdown of LPS by IAP, which was
        given in a diagram in one of the slides. The equations that arise from this metabolic
        pathway follow as:
        \begin{eqnarray*}
            \frac{dIAP}{dt} & = &  \alpha S_I - \rho_I IAP * LPS_G - \mu_I I\\
            \frac{dP}{dt} & = & \theta_I\left(\frac{LPS_G}{X_{ad} + LPS_G}\right) - \theta_R\frac{P}{P+K_R} \\
            \frac{dLPS_G}{dt} & = & \beta S_l - \rho_{L}IAP * LPS_G - P(LPS_G - LPS_B) \\
            \frac{dLPS_B}{dt} & = & DP(LPS_G - LPS_B) - \mu_L LPS_B
        \end{eqnarray*}
        After getting these equations we saw a variety of plots.
        The first plots were actual animal data taken from lab mice.
        We looked at the differences between a ``chow'' diet and a ``western'' diet
        over a multiple week period and how inflammation occurred in the animal.
        Further we saw how changes in diet could affect inflammation.
        Lastly, we were presented with a variety of numerical plots of data generated from the model
        comparing a healthy diet, a continued ``western'' diet a partially healthy diet and then those
        with stimulus where there was a stimulus to affect the production of IAP in a positive manner.

        \section{Nematocyst firing.}
        This problem I personally found more interesting than the rest, as we had a display of a
        single cell organism firing off an organelle, that we call a nematocyst, into another single
        cell organism, algae, for the purpose of capture.
        More generally nematocysts are employed by cnidarians (hydra, jellyfish, corals) and
        dinoflagellates (\emph{Polykrikos kofoiidi}), noting that the dinoflagellate mechanism
        seem to have more parts involved.
        Morphogensis of these organelles is mostly understood, but the mechanics of firing, on the other
        hand, is not.
        The firing is so fast that it is quite difficult to capture and analyze.

        The mathematical model that we employ is that of an accelerating rod that is transposed into
        the prey at velocity $V$. We have the variables $L$ as the length of the rod, $h$ as the distance
        between the prey and the rod, and $D$ as the diameter of the prey. The model follows from
        the non-dimensional NavierStokes equation
        \begin{gather*}
            \frac{\partial\mathbf{v}'}{\partial t'} +\mathbf{v}'\boldsymbol{\cdot}\nabla'\mathbf{v}' = -\nabla'p' + \frac{1}{\text{Re}}\nabla'^2\mathbf{v}'\\
            \text{Re}  = \frac{\rho LV}{\mu}
        \end{gather*}
        And, note that the Navier-Stokes equation can be put into this dimensionless form using the following
        definitions:
        \begin{align*}
            \mathbf{x}' &= \frac{\mathbf{x}}{L} & p' &= \frac{p}{\rho V^2} \\
            \mathbf{v}' &= \frac{\mathbf{v}}{V} & t' &= t\frac{V}{L} \\
            \nabla' &= L\nabla &
        \end{align*}
        where $L$ is the characteristic length and $V$ is the characteristic velocity.

        We then went over the code that would be written beginning with the mesh generation that we
        might create using finite element analysis or the finite difference method.
        Then the object would be to move the boundary around to accommodate for muscles, biochemistry,
        preferred curvature and other items.
        We notice that the model is fairly elastic with two different angles upon which we look at it:
        the first being a fiber based model where things hook together using beams and springs,
        the second being a continuum model using finite elements.

        We then considered a boundary layer approach where we model everything with spring mass dampers with a
        multitude of little springs to represent the boundary layer of the fluid.
        This lead to some interesting modeled results show as a video with a green background where the nematocysts
        never actually punctured the prey.
        We then looked at the difference across Reynolds numbers, namely for $\text{Re} = 0.02$ versus
        $\text{Re} = 200$.
        And we were mainly concerned with the ability of the nematocyst to reach the prey.
        Which, we plotted in various other plots.

        Ideas on future direction follow as: determine how nematocyst is ejected from the cell (applied
        force or pressure), and to study the mechanics of how the capsule is triggered.

        \section{Phage therapy.}
        So we didn't really get into the mathematical model behind the phage therapy, but we did get
        into what the basis of the problem is biologically, and it's quite interesting.
        The idea is to develop unique bacteriophages that are specifically tailored to the bacteria's
        transmembrane receptors.
        Note, bacteriophages essentially behave with a bacterial cell in the fashion that a virus
        would with a human cell.
        So the goal with the production of these bacteriophages is to administer them to patients
        with particularly officious bacterial infections.
        Furthermore, utilizing the technique for patients that are dealing with antiobiotic resistant
        bacteria or are particularly immunocompromised.

        We were very briefly shown two models one with delay differential equations (DDE) and one with
        ordinary differential equations (ODE).
        The ODE model was shown to fit the data better than the DDE model.
        Though, we can still gain insight from both models.
        For example, the DDE model tends to predict and earlier emergence of resistence as compared
        with the ODE model.
        However, in the long run the ODE model will be used for our project

        \section{Sickle cell pain episode prediction.}
        Similarly to the ``Phage therapy'' project we got an abbreviated presentation on the
        mathematical modeling of sickle cell anemia.
        The goal behind this project essentially is to see how pain works.
        We are using ODEs to model the cell dynamics of the hemoglobin distortions.
        The mechanism of pain is vascular occlusion and our goal is to properly model this.
        Further we're looking at inflammation generally associated with the anemia as well as
        how sleep affects recovery from an episode.

    \end{onehalfspace}
\end{document}
